\section{Related Work}
\label{sec:related}

LoRa capacity depends on airtime, radio configuration, timing, and
interference~\cite{bor2016scale,georgiou2017scale}. LoRaWAN studies also
distinguish received DATA from confirmed communication because ACK traffic
consumes constrained resources~\cite{vandenabeele2017scalability}.
Collision/capture experiments and ns-3 modeling motivate explicit reception
and interference treatment~\cite{rahmadhani2018concurrent,to2018lorans3}.
These studies do not calibrate the present simulator: inter-SF interference
and multi-interferer capture introduce further effects
\cite{croce2018orthogonality,heusse2023capacity}, and single-gateway results
do not establish multi-hop mesh performance.

Multi-hop LoRa research spans distinct relay/router architectures and
routing, access, and energy objectives
\cite{cotrim2020lorawanmesh,wong2024multihop}. Smart-Hop uses
spreading-factor-dependent range in a low-latency multi-hop design
\cite{ahmar2022smarthop}, while MLoRa targets economical coverage and
balanced energy use~\cite{feng2022cost}. Other approaches use SF clustering
to balance airtime~\cite{zhu2019improving} or time-on-air metrics in
distance-vector routing~\cite{pueyo2024minimalistic}. An ESP32 mesh library
and LoRaHop provide implementation evidence for proactive routing and
bidirectional LoRaWAN relaying, respectively
\cite{sole2022implementation,tian2023lorahop}. Their architectures and
hardware evidence cannot be treated as same-condition CPR baselines.

Source-route discovery and reuse are established by DSR
\cite{johnson1996dsr}, and finite ARQ exposes a delay--reliability trade-off
\cite{liu2004truncated}. Forward DATA and reverse ACK reception jointly
enter ETX's link cost~\cite{decouto2003etx}. Bursty low-power wireless
reception~\cite{srinivasan2008beta} further cautions against treating retry
outcomes as independent. CPR uses missing confirmation and local time
budgets; it does not implement an ETX estimator or derive reliability by
multiplying independent stage-success probabilities. Common-random-number
methods motivate shared-environment comparisons, although their variance
benefit is conditional on the system~\cite{glasserman1992crn}.

Redundant flooding creates contention and collision opportunities
\cite{ni1999broadcaststorm}. Trickle suppresses broadcasts after overhearing
consistent summaries~\cite{levis2004trickle}, while ExOR coordinates
opportunistic forwarding using acknowledgment information
\cite{biswas2005exor}. Official Meshtastic and MeshCore mechanisms likewise
include learned routes, retries, and flooding
\cite{meshtasticfirmware,meshcore}. Retry-to-flood recovery and ACK-assisted
suppression are therefore antecedents. CPR examines their specific
interaction with deadline admission, live-message confirmation, cache
evidence, and the point at which a relay send becomes committed.

The Meshtastic-like and MeshCore-like baselines are simplified behavioral
models. They omit official retry behavior; the former also omits next-hop
learning, while the latter uses RREQ/RREP rather than the documented
first-message DATA flood and delivery-report path learning. Comparisons
with these models cannot establish superiority over current product
firmware. Component controls isolate narrower CPR policy effects without
removing this external-fidelity limitation.
